\section{Other Holders and Earlier Cases}\label{app:cases}

\Cref{tab:cases,tab:history} code four contemporary and three historical cases with the rules of \cref{sec:coding}; this section gives the discussion that \cref{sec:cases} summarizes. Taiwan is outside the scope of \cref{sec:scope}, and we include it for what it suggests about control.

\begin{table}[p]
\centering
\scriptsize
\setlength{\tabcolsep}{2.5pt}
\caption{US demands on contemporary holders, coded through the model}\label{tab:cases}
\begin{tabularx}{\textwidth}{@{}>{\raggedright\arraybackslash}p{1.7cm}*{6}{>{\raggedright\arraybackslash}X}@{}}
\toprule
Holder, period & US demand & Instrument threatened & Compensation & Rival bid or counter-threat & Outcome & Reading in the model\\
\midrule
Netherlands (ASML), 2019--2026 & End EUV, then advanced and older immersion DUV and metrology sales to China; from 2026, servicing & Flow jurisdiction: zero de minimis (2023), SME FDP rule (2024), MATCH Act (proposed) & None; Supplement No.~4 status, a threat left unexercised & No Chinese offer for exit found; China is ASML's largest market (revenue, not an offer); Nexperia ban executed (2025) & Aligned on sales with lags; kept the licensing pen; servicing in China continued into 2026 & $Q\le0$ through jurisdiction (\cref{rem:firm}); hedge tolerated while no flow threat was credible (\cref{prop:counter}(b)); security channel largely untested\\
\addlinespace
Japan (Tokyo Electron), 2023--2026 & Control advanced equipment; in 2024, restrict sales and servicing in China & Flow jurisdiction: FDP rule not used but not ruled out & None found, including for the 2025--26 coercion; Supplement No.~4 status & 2024: ``severe economic retaliation'' threatened; 2025--26: 40 entities controlled and 40 watch-listed after a Taiwan statement, none of Japan's leading toolmakers among them & 23 equipment types licensed (2023), ``not aligning''; servicing continued into 2026 & Counter-threat raised the cost of the servicing concession (\cref{prop:counter}(a)); later punishment aimed at a security act\\
\addlinespace
South Korea (Samsung, SK hynix fabs in China), 2025--2026 & No expansion or upgrade of Chinese fabs & Flow jurisdiction: VEU status revoked (end-2025) & None; annual licences to operate (2026); November 2025 package silent on controls & None documented & Accepted; fabs operate under annual licence & Hedge policed by flow jurisdiction, rebuilt as a renewable status (\cref{prop:counter}(b))\\
\addlinespace
Germany (Zeiss, Trumpf), 2023--2024 & Join equipment controls; ``primary target is Zeiss'' & None documented & None & None documented & Reportedly declined; slower licensing of laser exports to China & Demand of low marginal value while Dutch licences govern the tools that carry Zeiss optics (\cref{cor:rationing})\\
\addlinespace
Taiwan (TSMC) and US-sited assets, 2020--2026 & Relocate leading-edge capacity; ``50-50'' floated (2025) & Tariffs; Section 232 preference for firms investing in the US & CHIPS awards; tariff preference tied to US investment & China a threat, not a bidder & TSMC's US commitment from about \$12bn to \$165bn; US control rights only over US-sited assets & Control sought where $\Delta_c$ is high (\cref{cor:control}(c)); outside scope\\
\bottomrule
\end{tabularx}
\par\smallskip
\begin{minipage}{\textwidth}\scriptsize
\emph{Sources by row.} Netherlands: \cref{sec:background} and \citet{bis2024fdpr, asml2026annual, zakaria2026match, baumgartner2026match, kst32637720, kst32637_717}. Japan: \citet{allen2023translation, hoganlovells2023japan, taipeitimes2024retaliation, kokkai2025takaichi, mofcom2026japan11, mofcom2026japan27, teer2026dismantling, zakaria2026match}. South Korea: \citet{bis2025veu, afs2025veu, skhynix2026prospectus, reuters2025koreanlicence, census2026fablicense, whitehouse2025korea}. Germany: \citet{barkin2024zeiss, burger2023customs, kahle2023trumpf}. Taiwan and US-sited assets: \citet{younger2025fiftyfifty, focustaiwan2025fiftyfifty, tsmc2020arizona, tsmc2025expand, proc11002, tsmc2026annual, mp2025dod, ussteel2025nsa, intel2025agreement}. ``None documented'' and ``None found'' mean that the sources we consulted record none; they cannot exclude private linkage or compensation.
\end{minipage}
\end{table}

\begin{table}[p]
\centering
\scriptsize
\setlength{\tabcolsep}{2.5pt}
\caption{Historical cases, coded through the model}\label{tab:history}
\begin{tabularx}{\textwidth}{@{}>{\raggedright\arraybackslash}p{1.7cm}*{6}{>{\raggedright\arraybackslash}X}@{}}
\toprule
Holder, period & US demand & Instrument threatened & Compensation & Rival bid or counter-threat & Outcome & Reading in the model\\
\midrule
West Germany, 1961--1967 (offsets) & Offset the foreign-exchange cost of US troops & Conditional linkage: troop levels, explicit in 1966 & Holder paid: arms purchases, \$500m in US securities, gold pledge & German option to convert dollars into gold, given up & Redeployment of up to 35,000 US personnel agreed (1967); Erhard fell (1966) & Linkage extracting sunk payments; repeated renegotiation (\cref{prop:slack}(d))\\
\addlinespace
United Kingdom, 1957 (CoCom China differential) & Keep stricter controls on China than on the Soviet bloc & None; Washington chose no sanctions & None & China trade; 12 of 16 members for abolition & Differential abolished; US kept its own embargo & Secondary demand left without linkage (\cref{cor:rationing})\\
\addlinespace
West Germany, 1975--1977 (Brazil nuclear deal) & Defer or stop enrichment and reprocessing transfers & High-level pressure; no explicit linkage, though Bonn feared for the US commitment & Offered to Brazil (fuel guarantee), not to Bonn & Brazil as buyer; Bonn cited its contract and the NPT & Contract kept; informal restraint on future transfers by 1978 & Concession on the forward margin; sunk deal kept\\
\bottomrule
\end{tabularx}
\par\smallskip
\begin{minipage}{\textwidth}\scriptsize
\emph{Sources by row.} Offsets: \citet{frus1995foreignecon, frus1995westerneurope, frus1999germany, bundesbank1967blessing}. United Kingdom: \citet{frus1989economicdefense, hansard1957china}. Brazil deal: \citet{frus2015southamerica, frus2014westerneurope, frus2015nonproliferation}.
\end{minipage}
\end{table}

\paragraph{What the comparison suggests.} We found no holder paid in cash for aligning on a chokepoint. The currency is status that BIS grants and withdraws by rule (Supplement No.~4, validated end-user status, annual licences), which we code as exit costs. Every holder in scope has, we think, a nonpositive premium, so the tables cannot test the threshold of \cref{cor:master}. They speak to the instruments Washington uses and to where holders resist. Where money moved, it moved to the patron (the offsets) or to a third party, as when Washington offered Brazil, not Bonn, a fuel guarantee in 1977 \citep{frus2015nonproliferation}.

We read Korea as flow jurisdiction rebuilt as a renewable status. BIS moved the Samsung and SK hynix fabs in China from validated end-user status to annual licences that allow operation but not expansion or upgrades \citep{afs2025veu, skhynix2026prospectus, reuters2025koreanlicence}, and the trade and security package of November 2025 did not mention export controls \citep{whitehouse2025korea}.

The Japanese case is consistent with Chinese punishment aimed at a security act. After Prime Minister Takaichi told the Diet on 7 November 2025 that a Taiwan contingency could become a survival-threatening situation for Japan \citep{kokkai2025takaichi}, China placed forty Japanese entities on its export-control list and forty more on its watch list \citep{mofcom2026japan11, mofcom2026japan27}. None is among Japan's leading makers of lithography, deposition and etch tools or photoresists (Tokyo Electron, Nikon, Canon, JSR, Tokyo Ohka, Shin-Etsu). Beijing spared the holders whose inputs it needs. We found no US compensation \citep{teer2026dismantling}, which \cref{cor:master} leads us to expect while Japan's premium stays negative. Coordination among holders runs through Washington: Japan and South Korea have signed Pax Silica, the US-led initiative on AI chip supply chains and export controls, and TNW reported in June 2026 that the EU was set to join \citep{tnw2026paxsilica}.

Taiwan is consistent with what the model expects of a patron that can move production. TSMC's US commitment rose from about USD 12 billion in 2020 to USD 165 billion in 2025 \citep{tsmc2020arizona, tsmc2025expand}, and the control rights Washington has taken all attach to assets on US soil (\cref{sec:preventive}). We found no US demand for equity in a foreign chipmaking-equipment firm. A demand for a controlling vote in ASML presupposes that relocating EUV production and servicing is infeasible.

Resistance clusters on the installed base and on signed deals: Tokyo and The Hague kept servicing Chinese fabs, and Bonn kept its contract with Brazil. In the model these concessions cost the holder much relative to the remaining bite, but a cost minimizer facing expensive enforcement would also tolerate them.

Explicit links across issues appear in Johnson's tie between troops and offsets in 1966, in the posture cut in Germany and the threat to cut off trade with Spain after it denied base access in 2026 \citep{martinez2026germany, roque2026spain, omullony2026ankara}, and in the BIS rule of July 2026 that cited the UAE's role in Operation Epic Fury when it granted licence-free chip access \citep{bis2026uae}. None concerns a European holder's chip controls.

\paragraph{Where the comparison strains the model.} The offset payments were sunk concessions, which escape the slack cap (\cref{prop:slack}(b)), and renegotiation recurred from 1961 to 1967 \citep{frus1995foreignecon, frus1999germany, frus1995westerneurope}, as \cref{prop:slack}(d) leads us to expect of a patron whose stakes are small relative to its later demands. But a complete-information patron does not both collect and cut, and in 1967 Washington agreed to redeploy up to 35,000 personnel. If the troops were above Washington's ex post optimum, as balance-of-payments pressure and the Mansfield resolution suggest, keeping them was a transfer and the offsets were its price \citep{gavin2004gold, zimmermann2002money}. If not, the cut is unexplained, and FRUS does not let us choose. In 1957 Washington declined to back its China differential with linkage, as \cref{cor:rationing} implies for a demand worth little, but twelve of sixteen CoCom members backed abolition \citep{frus1989economicdefense, mastanduno1992economic}, a collective exit our single-holder model does not represent. In the Brazil case Bonn feared for a US commitment that Washington did not tie to its pressure \citep{frus2014westerneurope}, so conditionality existed only in the holder's beliefs.
